% GENERATED by tools/coverage.py from the lesson plans -- do not edit by hand.
\subsection{Objectives per goal and week (generated)}
Each number is how many of that week's \enquote{The student can\ldots} objectives
serve the goal. Blank = not a focus that week (it may still be practiced).

{\small\setlength{\tabcolsep}{3pt}
\begin{longtable}{@{}lcccccccccccc@{}}
\toprule
Goal & 1 & 2 & 3 & 4 & 5 & 6 & 7 & 8 & 9 & 10 & B1 & B2\\ \midrule\endhead
J3.G1 & 2 &  &  &  & 1 &  & 1 & 1 & 1 &  &  & 1\\
J3.G2 &  & 3 &  &  & 1 &  &  &  & 1 & 1 & 2 & 1\\
J3.G3 &  &  & 3 & 1 &  &  &  &  &  &  & 1 & \\
J3.G4 &  &  & 1 & 3 &  &  &  &  &  &  &  & \\
J3.G5 &  &  &  &  & 1 & 2 & 2 & 3 & 1 & 2 &  & 1\\
J3.G6 & 1 & 2 & 1 & 1 & 1 & 1 & 2 & 1 & 1 & 2 & 1 & 1\\
J3.G7 & 1 &  & 1 & 1 & 1 & 1 & 1 & 1 & 1 & 1 &  & 1\\
\bottomrule\end{longtable}}

\subsection{All lesson objectives (generated)}
\begin{longtable}{@{}p{1.5cm}p{0.72\linewidth}p{1.6cm}@{}}
\toprule Week & Objective & Goals\\ \midrule\endhead
\multicolumn{3}{@{}l}{\textbf{1: Patterns in Patterns}}\\
1.1 & The student can clap a pattern inside a pattern and build a repeat inside a repeat in Scratch that draws a window of four squares. & J3.G1\\
1.2 & The student can use a \enquote{square} My Block inside a repeat to draw a flower of six squares, reading the turn from the turn card. & J3.G1\\
1.3 & The student can find and fix an inner repeat that sits under the outer repeat instead of inside it. & J3.G6\\
1.4 & The student can make a pattern of their own and say whose idea it came from when they borrow one. & J3.G7\\
\midrule
\multicolumn{3}{@{}l}{\textbf{2: Plan It on Paper}}\\
2.1 & The student can split a party, and then a Scratch program, into named parts on a plan card, with a short step list for each part. & J3.G2\\
2.2 & The student can build a program one part at a time and test each part on its own by clicking its script. & J3.G2, J3.G6\\
2.3 & The student can find and fix a part that works alone but not with the others because a message has the wrong name. & J3.G6\\
2.4 & The student can share out the parts of a plan so that everyone on the team has a job. & J3.G2\\
\midrule
\multicolumn{3}{@{}l}{\textbf{3: Data Detectives}}\\
3.1 & The student can run a family survey, keep count with tally marks, and draw a bar graph that shows which fruit won. & J3.G3\\
3.2 & The student can build a Scratch survey that asks a question, checks the answer with if/else and counts it in a variable. & J3.G3\\
3.3 & The student can find one answer counted in two places, such as \enquote{Apple} and \enquote{apple}, and fix the tally and the graph. & J3.G3, J3.G6\\
3.4 & The student can say why we only ask what we need, never ask for personal information, and why asking only one kind of person gives a graph that is not fair. & J3.G4, J3.G7\\
\midrule
\multicolumn{3}{@{}l}{\textbf{4: How Do Computers Learn?}}\\
4.1 & The student can teach a partner to sort animal cards from examples and say which features, such as legs, wings or fur, helped. & J3.G4\\
4.2 & The student can test the \enquote{robot} with new cards, tally its right and wrong predictions, and show them in a bar graph. & J3.G3, J3.G4\\
4.3 & The student can explain how too few or one-sided examples lead to wrong guesses, and fix it by adding different examples. & J3.G4, J3.G6\\
4.4 & The student can say that AI can be wrong and that we never tell a chatbot personal information. & J3.G7\\
\midrule
\multicolumn{3}{@{}l}{\textbf{5: Play Day: Blocks to Words}}\\
5.1 & The student can match Scratch blocks to their Python word twins, such as \enquote{move 100 steps} to \texttt{t.forward(100)} and \enquote{repeat 4} to \texttt{for i in range(4):}. & J3.G5\\
5.2 & The student can find and fix the bug in a nested repeat, a program built in two parts and a counter, by testing and changing one thing at a time. & J3.G6\\
5.3 & The student can plan a pattern in three named parts and build it with a repeat inside a repeat. & J3.G1, J3.G2\\
5.4 & The student can say one thing they are proud of in their own work and give credit for an idea they borrowed. & J3.G7\\
\midrule
\multicolumn{3}{@{}l}{\textbf{6: Hello, Python!}}\\
6.1 & The student can open a starter program in Thonny, press Run, point to the Editor, the Shell and the turtle window, and change a number to make the turtle walk further. & J3.G5\\
6.2 & The student can match \texttt{t.forward(100)} and \texttt{t.left(90)} to the Scratch move and turn blocks, and type one new line to finish a square. & J3.G5\\
6.3 & The student can find a typo such as \texttt{forwad} or a missing bracket by reading the line number and the last line of Thonny's red error text, and fix it. & J3.G6\\
6.4 & The student can save a program with a clear file name that uses a nickname, never a real name. & J3.G7\\
\midrule
\multicolumn{3}{@{}l}{\textbf{7: Repeat in Python}}\\
7.1 & The student can read \texttt{for i in range(4):} as \enquote{repeat 4} and change its number to draw a square, or a triangle with the number from the turn card. & J3.G5\\
7.2 & The student can find a line that is not pushed in, explain that it runs only once, and fix it by pushing it in with the Tab key. & J3.G6\\
7.3 & The student can point to the inner and the outer repeat of a repeat inside a repeat in Python, and change the outer number to draw a flower of six squares. & J3.G1, J3.G5\\
7.4 & The student can ask for help well by saying what they expected and what happened, and asks a grown-up they know instead of a website or a chatbot. & J3.G6, J3.G7\\
\midrule
\multicolumn{3}{@{}l}{\textbf{8: Colors and My Own Pattern}}\\
8.1 & The student can change a color name and a pen size in a working Python program, one change at a time, and run it after each change. & J3.G5\\
8.2 & The student can find a misspelled color name with the help of Thonny's red \enquote{bad color string} text and fix the spelling. & J3.G6\\
8.3 & The student can remix a working flower into a two-color pattern and write a credit line that says whose program it started from, using only a nickname. & J3.G5, J3.G7\\
8.4 & The student can say that \texttt{def square():} is like a Scratch My Block called square and use it inside a repeat (stretch). & J3.G1, J3.G5\\
\midrule
\multicolumn{3}{@{}l}{\textbf{9: Project Studio}}\\
9.1 & The student can plan their Big Build on the plan card: choose Scratch or Python, draw it, and name its parts and steps before building. & J3.G2\\
9.2 & The student can build their project one part at a time and test after each part: a Scratch game or art app with a nested repeat or a My Block and a score or counter, or a Python postcard with a \texttt{for} loop and two colors. & J3.G1, J3.G5\\
9.3 & The student can check a partner's project against the plan card, and find and fix a bug in blocks or a typo in Python. & J3.G6\\
9.4 & The student can choose colors and sizes that everyone can see, and write down whose idea they used. & J3.G7\\
\midrule
\multicolumn{3}{@{}l}{\textbf{10: Graduation Showcase}}\\
10.1 & The student can show their Big Build, say how it matches their plan card, and explain what one part of their code does in their own words. & J3.G2, J3.G5\\
10.2 & The student can run a Python program in Thonny, fix a typo by reading the red error message, and change a number and a color to make it their own. & J3.G5, J3.G6\\
10.3 & The student can test their project in a last-minute bug check and fix one thing without breaking the rest. & J3.G6\\
10.4 & The student can say what to do if a website asks for their name, and give kind feedback and credit at the gallery walk. & J3.G7\\
\midrule
\multicolumn{3}{@{}l}{\textbf{B1: Robot Olympics}}\\
B1.1 & The student can sort a row of number cards by following a written step rule and say how they know the row is sorted. & J3.G2\\
B1.2 & The student can solve a paper maze by hand with the \enquote{hand on the right wall} rule and write the moves as a step list a partner can follow. & J3.G2\\
B1.3 & The student can find and fix the bug in a sorting rule or a maze step list by acting it out. & J3.G6\\
B1.4 & The student can tally the swaps in two sorting races and compare them in a bar graph. & J3.G3\\
\midrule
\multicolumn{3}{@{}l}{\textbf{B2: Free Build}}\\
B2.1 & The student can choose Scratch or Python and plan a small project in two or three named parts before building. & J3.G2\\
B2.2 & The student can build the project part by part, using a repeat inside a repeat or a My Block in Scratch, or changing numbers, colors and \texttt{for} loops in Python. & J3.G1, J3.G5\\
B2.3 & The student can test a partner's project, describe a bug as \enquote{I expected \ldots, but it \ldots}, and fix a bug in their own project. & J3.G6\\
B2.4 & The student can give credit in their project for an idea they remixed. & J3.G7\\
\midrule
\bottomrule\end{longtable}
